\chapter*{Conclusion G\'en\'erale et Perspectives}
\addcontentsline{toc}{chapter}{Conclusion G\'en\'erale et Perspectives}
\markboth{Conclusion G\'en\'erale et Perspectives}{Conclusion G\'en\'erale et Perspectives}

\section*{1. Synth\`ese des Travaux R\'ealis\'es}
\par Ce projet de fin d'\'etudes d'ing\'enieur, r\'ealis\'e au sein de l'agence num\'erique \textbf{Designet Web Agency}, a port\'e sur la conception et le d\'eveloppement d'une plateforme web semi-agentique de reconnaissance de surface d'attaque, associant une corr\'elation par graphe de connaissances, une g\'en\'eration de rem\'ediations assist\'ee par intelligence artificielle locale contrainte et une cha\^ine de livraison DevSecOps continue.
\par Pour surmonter les limites des d\'emarches manuelles traditionnelles, nous avons d\'evelopp\'e une architecture en microservices d\'ecoupl\'ee et r\'esiliente. Le socle applicatif comprend un back-end en \textbf{Laravel 11}, cinq microservices en \textbf{Python Flask} (reconnaissance multi-outils, s\'ecurit\'e offensive / sandbox, OSINT passif, IA d'assistance, passerelle API), un moteur d'ex\'ecution \'ev\'enementiel pilot\'e par \textbf{Redis Streams}, une couche de donn\'ees unifi\'ee sous \textbf{PostgreSQL 16 / Apache AGE}, et un moteur d'inf\'erence locale \textbf{Ollama} h\'ebergeant le mod\`ele d\'edi\'e au code \texttt{qwen2.5-coder}.

\section*{2. Bilan des Contributions et R\'esultats}
\par Les travaux conduits ont abouti \`a des r\'esultats scientifiques et op\'erationnels majeurs :
\begin{itemize}
    \item \textbf{Orchestration unifi\'ee des scanners :} Encapsulation homog\`ene de Nmap, Nuclei, Gobuster, Subfinder et WPScan avec gestion de profils d'intensit\'e et injection de jitter anti-d\'etection.
    \item \textbf{Mod\'elisation en graphe et calcul de blast radius :} Projection relationnelle et openCypher de la surface d'attaque avec calcul instantan\'e du rayon d'impact via l'algorithme BFS sous NetworkX ($<$75 ms).
    \item \textbf{Assistance IA locale et Remediation-as-Code tra\c{c}able :} Production de scripts de durcissement (Bash, Ansible, Dockerfile) au format JSON contraint avec citations directes vers les journaux bruts, garantissant une confidentialit\'e 100\,\% locale.
    \item \textbf{S\'ecurit\'e intrins\`eque DevSecOps :} D\'efense en profondeur (conteneurs non-root, socket proxy) et pipeline CI/CD automatis\'e (SAST, SCA, SBOM CycloneDX, scans Trivy et signature Cosign).
    \item \textbf{Validation exp\'erimentale rigoureuse :} Validation de 140/140 tests automatis\'es (PHPUnit et Pytest) et tenue sous charge de 50 utilisateurs simultan\'es sous Locust sur l'instance de production \url{https://aymenazizi.dijaly.com/}.
\end{itemize}

\section*{3. R\'eflexion Personnelle et Comp\'etences Acquises}
\par Ce projet de fin d'\'etudes a constitu\'e une exp\'erience professionnelle et humaine particuli\`erement enrichissante. Il m'a permis de consolider une double comp\'etence de pointe :
\begin{itemize}
    \item \textbf{Comp\'etences Techniques :} Ma\^itrise avanc\'ee des architectures distribu\'ees en microservices, ing\'enierie des syst\`emes asynchrones \'ev\'enementiels (Redis Streams), int\'egration de mod\`eles LLM dans des flux contraints, requ\^etage openCypher sur bases de graphes hybrides et automatisation DevSecOps de bout en bout.
    \item \textbf{Comp\'etences Humaines et M\'ethodologiques (\textit{Soft Skills}) :} Autonomie et rigueur dans la r\'esolution de probl\`emes d'ing\'enierie complexes, pilotage it\'eratif selon le cadre Agile Scrum, communication technique efficace lors des revues de sprint avec le CTO, et gestion rigoureuse des priorit\'es et du temps.
\end{itemize}

\section*{4. Perspectives d'\'Evolution}
\par Plusieurs perspectives d'\'evolution sont envisag\'ees pour \'etendre les capacit\'es de la plateforme :
\begin{itemize}
    \item \textbf{\`A court terme :} Int\'egration de scanners compl\'ementaires (Nikto, OpenVAS, analyseurs d'Infrastructure as Code tels que tfsec ou Checkov).
    \item \textbf{\`A moyen terme :} D\'eveloppement d'une application mobile d\'edi\'ee aux analystes SOC pour la notification instantan\'ee des alertes critiques et le suivi mobile des audits.
    \item \textbf{\`A long terme :} R\'ealisation d'un entra\^inement sp\'ecifique (\textit{fine-tuning}) de mod\`eles LLM compacts sur des corpus de vuln\'erabilit\'es industrielles et enrichissement du moteur de corr\'elation par apprentissage automatique pour la d\'etection pr\'edictive des vecteurs d'attaque.
\end{itemize}
